% Einheit 3 von 8 – Pfadregeln bei unabhängigen Stufen (bank/zufallsexperimente-und-pfadregeln/e3.jsonl, zusammenbau v0.1)
%% TODO Zeitmarke Einheit 3: Marken-Zeile nicht lesbar
%% TODO Zweigzeile Teil 1: was hier gelernt wird (Satz in Schülersprache aus dem Einheitstitel) – Einheit 3
\einheitenkopf[e3]{Einheit 3 von 8 · Pfadregeln bei unabhängigen Stufen}
\zweigzeile{Abitur GK}
\setcounter{aufgabe}{16}

%% TODO Titel: Ich-kann-Satz fehlt in der Bank (Platzhalter: Kettenname) – Nr. 17
%% TODO Anweisung: ein Satz über der ersten Teilaufgabe fehlt in der Bank – Nr. 17
\begin{aufgabe}{Pfadregeln bei unabhängigen Stufen}
%% TODO \rechenplatz{n}: Schrittzahl des Grundfalls fehlt in der Bank (nur bei fester Schreibform, 2.3 b)
\begin{teile}
\teil Ein Glücksrad wird dreimal gedreht. Bleibt P(rot) auf jeder Stufe gleich?\\ \kreuz{p bleibt gleich}\\ \kreuz{p ändert sich}
\teil Ein Glücksrad zeigt rot (R) mit 30\,\%, sonst blau (B); es wird zweimal gedreht. Beschrifte den Baum. P(erst rot, dann blau)?

\baumzwei{R/,B/}{R/,B/}{R/,B/}
\teil Ein Glücksrad liefert einen Treffer mit $0{,}4$. Es wird fünfmal gedreht. P(erst drei Nieten, dann zwei Treffer)? P = \leerfeld
\teil Eine Ampel ist bei Ankunft mit $0{,}35$ grün; vier Ampeln schalten unabhängig. P(alle vier grün)? P = \leerfeld
\teil Eine Urne enthält 3 rote und 5 blaue Kugeln; zweimal mit Zurücklegen. P(beide gleichfarbig)? P = \leerfeld
\teil Ein Glücksrad zeigt Gewinn (G) mit $0{,}3$, sonst Niete (N); zweimal. Vollständiger Baum. A: genau einmal Gewinn. B: mindestens einmal Gewinn. P(A), P(B), Gegenereignis zu B in Worten?

\baumzwei{G/,N/}{G/,N/}{G/,N/}
\teil Ein Versuch gelingt mit $0{,}15$. P(mindestens ein Erfolg in sechs Versuchen)? P = \leerfeld
\teil Lena würfelt höchstens dreimal und hört bei der ersten Sechs auf. P(sie würfelt eine Sechs)? P = \leerfeld
\teil A und B würfeln abwechselnd, A beginnt; wer zuerst eine Sechs hat, gewinnt; jeder würfelt höchstens zweimal. P(A gewinnt), P(B gewinnt)? P(A) = \leerfeld, P(B) = \leerfeld
\teil Würfel C trägt 2, 2, 2, 6, 6, 6, Würfel D trägt 5, 5, 5, 1, 1, 1. Je Runde wird jeder einmal geworfen; zwei Runden. P(in beiden Runden zeigt C mehr als D)? \hfill (Abitur 2026 LK) \\ P = \leerfeld
\end{teile}
\end{aufgabe}

%% TODO Titel: Ich-kann-Satz fehlt in der Bank (Platzhalter: Kettenname) – Nr. 18
%% TODO Anweisung: ein Satz über der ersten Teilaufgabe fehlt in der Bank – Nr. 18
\begin{aufgabe}{Gleich wahrscheinliche Ergebnisse eines zweistufigen Experiments über die Pfadwahrscheinlichkeiten ermitteln}
\begin{teile}
\teil Ein Würfel trägt zweimal X und viermal Y; ein Glücksrad zeigt rot mit $\frac{1}{2}$, blau und weiß mit je $\frac{1}{4}$. Beide werden benutzt. Welche drei der sechs Ergebnisse sind gleich wahrscheinlich?
\end{teile}
\end{aufgabe}

%% TODO Titel: Ich-kann-Satz fehlt in der Bank (Platzhalter: Kettenname) – Nr. 19
%% TODO Anweisung: ein Satz über der ersten Teilaufgabe fehlt in der Bank – Nr. 19
\begin{aufgabe}{Kosten je Erfolg zweier Varianten über Pfadwahrscheinlichkeiten vergleichen}
\begin{teile}
\teil Blumenzwiebel A kostet $0{,}40$\,€, treibt mit 90\,\% aus und blüht dann mit 80\,\%; Zwiebel B kostet $0{,}30$\,€, treibt mit 75\,\% aus und blüht dann mit 70\,\%. Welche ist je blühender Pflanze günstiger?
\end{teile}
\end{aufgabe}

%% TODO Titel: Ich-kann-Satz fehlt in der Bank (Platzhalter: Kettenname) – Nr. 20
%% TODO Anweisung: ein Satz über der ersten Teilaufgabe fehlt in der Bank – Nr. 20
\begin{aufgabe}{Pfadregeln bei unabhängigen Stufen – Fehler finden, Begründen, Anwendung, Darstellungswechsel}
\begin{teile}
\teil P(Treffer) $= 0{,}3$, drei Versuche. Ole rechnet P(dreimal Treffer) $= 3 \cdot 0{,}3 = 0{,}9$. Finde den Fehler.
\teil Begründe: Beim Ziehen mit Zurücklegen stehen auf jeder Stufe des Baums dieselben Zahlen.
\teil Ein Rauchmelder schlägt bei Rauch mit $0{,}95$ an. Im Flur hängen zwei unabhängige Melder. P(mindestens einer schlägt an)? P = \leerfeld
\teil Schreibe zum Baum den Term für P(genau einmal T) und berechne ihn.

\baumzwei{T/0{,}6,N/0{,}4}{T/0{,}6,N/0{,}4}{T/0{,}6,N/0{,}4}
\end{teile}
\end{aufgabe}
